\chapter{2013年大纲卷（理）}

\section{单选题}

\begin{problem}
设集合 \(A = \{1, 2, 3\}\)，\(B = \{4, 5\}\)，\(M = \{x \mid x = a + b, a \in A, b \in B\}\)，则 \(M\) 中元素的个数为（\quad）

\choices
    {\(3\)}
    {\(4\)}
    {\(5\)}
    {\(6\)}
\end{problem}

\begin{answer}
B
\end{answer}

\begin{solution}
由题意，\(a+b\) 的所有可能值为 \(1+4\)，\(1+5\)，\(2+4\)，\(2+5\)，\(3+4\)，\(3+5\)，所以 \(M=\{5,6,7,8\}\).因此 \(M\) 中有 \(4\) 个不同元素，故选 B.
\end{solution}

\begin{problem}
\(\left(1 + \sqrt{3}\i\right)^3 =\)（\quad）

\choices
    {\(-8\)}
    {\(8\)}
    {\(-8\i\)}
    {\(8\i\)}
\end{problem}

\begin{answer}
A
\end{answer}

\begin{solution}
因为 \(1+\sqrt{3}\i=2\left(\cos\frac{\pi}{3}+\i\sin\frac{\pi}{3}\right)\)，所以由棣莫弗定理，\(\left(1+\sqrt{3}\i\right)^3=2^3\left(\cos\pi+\i\sin\pi\right)=-8\).故选 A.
\end{solution}

\begin{problem}
已知向量 \(\bs{m} = (\lambda + 1, 1)\)，\(\bs{n} = (\lambda + 2, 2)\)，若 \((\bs{m} + \bs{n}) \perp (\bs{m} - \bs{n})\)，则 \(\lambda =\)（\quad）

\choices
    {\(-4\)}
    {\(-3\)}
    {\(-2\)}
    {\(-1\)}
\end{problem}

\begin{answer}
B
\end{answer}

\begin{solution}
由题设，\(\bs{m}+\bs{n}=(2\lambda+3,3)\)，\(\bs{m}-\bs{n}=(-1,-1)\).两向量垂直等价于内积为零，因此 \((2\lambda+3)(-1)+3(-1)=0\)，即 \(-2\lambda-6=0\)，解得 \(\lambda=-3\).故选 B.
\end{solution}

\begin{problem}
已知函数 \(f(x)\) 的定义域为 \((-1,0)\)，则函数 \(f(2x+1)\) 的定义域为（\quad）

\choices
    {\((-1,1)\)}
    {\(\left(-1,-\frac{1}{2}\right)\)}
    {\((-1,0)\)}
    {\(\left(\frac{1}{2},1\right)\)}
\end{problem}

\begin{answer}
B
\end{answer}

\begin{solution}
要使 \(f(2x+1)\) 有意义，必须有 \(-1<2x+1<0\).两边同时减去 \(1\)，再除以正数 \(2\)，得 \(-1<x<-\frac{1}{2}\).因此定义域为 \(\left(-1,-\frac{1}{2}\right)\)，故选 B.
\end{solution}

\begin{problem}
函数 \(f(x)=\log_2\left(1+\frac{1}{x}\right)\)（\(x>0\)）的反函数 \(f^{-1}(x)=\)（\quad）

\choices
    {\(\frac{1}{2^x-1}\)（\(x>0\)）}
    {\(\frac{1}{2^x-1}\)（\(x\neq0\)）}
    {\(2^x-1\)（\(x\in\R\)）}
    {\(2^x-1\)（\(x>0\)）}
\end{problem}

\begin{answer}
A
\end{answer}

\begin{solution}
令 \(y=f(x)\)，则 \(2^y=1+\frac{1}{x}\)，从而 \(x=\frac{1}{2^y-1}\).当 \(x>0\) 时，\(1+\frac{1}{x}>1\)，故 \(y>0\).交换变量名，得 \(f^{-1}(x)=\frac{1}{2^x-1}\)（\(x>0\)）.故选 A.
\end{solution}

\begin{problem}
已知数列 \(\{a_n\}\) 满足 \(3a_{n+1}+a_n=0\)，\(a_2=-\frac{4}{3}\)，则 \(\{a_n\}\) 的前 10 项和等于（\quad）

\choices
    {\(-6(1-3^{-10})\)}
    {\(\frac{1}{9}(1-3^{10})\)}
    {\(3(1-3^{-10})\)}
    {\(3(1+3^{-10})\)}
\end{problem}

\begin{answer}
C
\end{answer}

\begin{solution}
由递推关系得 \(a_{n+1}=-\frac{1}{3}a_n\).又由 \(a_2=-\frac{1}{3}a_1=-\frac{4}{3}\)，得 \(a_1=4\).所以 \(\{a_n\}\) 是首项为 \(4\)、公比为 \(-\frac{1}{3}\) 的等比数列，从而 \(S_{10}=4\frac{1-\left(-\frac{1}{3}\right)^{10}}{1-\left(-\frac{1}{3}\right)}=3(1-3^{-10})\).故选 C.
\end{solution}

\begin{problem}
\((1+x)^8(1+y)^4\) 的展开式中 \(x^2y^2\) 的系数是（\quad）

\choices
    {\(56\)}
    {\(84\)}
    {\(112\)}
    {\(168\)}
\end{problem}

\begin{answer}
D
\end{answer}

\begin{solution}
要得到 \(x^2y^2\)，必须分别从两个因式中取出 \(x^2\) 和 \(y^2\).由二项式定理，该项系数为 \(\mathrm C_8^2\mathrm C_4^2=28\times6=168\).故选 D.
\end{solution}

\begin{problem}
椭圆 \(C:\frac{x^2}{4}+\frac{y^2}{3}=1\) 的左、右顶点分别为 \(A_1\)，\(A_2\)，点 \(P\) 在 \(C\) 上且直线 \(PA_2\) 斜率的取值范围是 \([-2,-1]\)，那么直线 \(PA_1\) 斜率的取值范围是（\quad）

\choices
    {\(\left[\frac{1}{2},\frac{3}{4}\right]\)}
    {\(\left[\frac{3}{8},\frac{3}{4}\right]\)}
    {\(\left[\frac{1}{2},1\right]\)}
    {\(\left[\frac{3}{4},1\right]\)}
\end{problem}

\begin{answer}
B
\end{answer}

\begin{solution}
椭圆的左、右顶点为 \(A_1=(-2,0)\)，\(A_2=(2,0)\).设 \(P=(x,y)\)，直线 \(PA_1\)，\(PA_2\) 的斜率分别为 \(k_1\)，\(k_2\).由椭圆方程，\(y^2=3\left(1-\frac{x^2}{4}\right)=\frac{3}{4}(4-x^2)\)，因此 \(k_1k_2=\frac{y^2}{(x+2)(x-2)}=-\frac{3}{4}\).已知 \(k_2\in[-2,-1]\)，故 \(k_1=-\frac{3}{4k_2}\in\left[\frac{3}{8},\frac{3}{4}\right]\).故选 B.
\end{solution}

\begin{problem}
若函数 \(f(x)=x^2+ax+\frac{1}{x}\) 在 \(\left(\frac{1}{2},+\infty\right)\) 是增函数，则 \(a\) 的取值范围是（\quad）

\choices
    {\([-1,0]\)}
    {\([-1,+\infty)\)}
    {\([0,3]\)}
    {\([3,+\infty)\)}
\end{problem}

\begin{answer}
D
\end{answer}

\begin{solution}
在定义域上 \(f'(x)=2x+a-\frac{1}{x^2}\).令 \(g(x)=\frac{1}{x^2}-2x\)，则函数递增的必要条件是对任意 \(x>\frac{1}{2}\) 有 \(a\geqslant g(x)\).又 \(g'(x)=-\frac{2}{x^3}-2<0\)，所以 \(g\) 在 \(\left(\frac{1}{2},+\infty\right)\) 上递减，且 \(g(x)\) 的上确界为 \(\lim_{x\to(1/2)^+}g(x)=3\).因此必要条件是 \(a\geqslant3\).反之，当 \(a\geqslant3\) 时，\(g(x)<3\leqslant a\)，所以 \(f'(x)>0\)，函数确实递增.故 \(a\in[3,+\infty)\)，选 D.
\end{solution}

\begin{problem}
已知正四棱柱 \(ABCD-A_1B_1C_1D_1\) 中，\(AA_1=2AB\)，则 \(CD\) 与平面 \(BDC_1\) 所成角的正弦值等于（\quad）

\choices
    {\(\frac{2}{3}\)}
    {\(\frac{\sqrt{3}}{3}\)}
    {\(\frac{\sqrt{2}}{3}\)}
    {\(\frac{1}{3}\)}
\end{problem}

\begin{answer}
A
\end{answer}

\begin{solution}
设正方形底面的边长为 \(s\)，建立直角坐标系，取 \(A=(0,0,0)\)，\(B=(s,0,0)\)，\(D=(0,s,0)\)，\(C=(s,s,0)\)，\(C_1=(s,s,2s)\).平面 \(BDC_1\) 的两个方向向量可取 \(\overrightarrow{BD}=(-s,s,0)\) 和 \(\overrightarrow{BC_1}=(0,s,2s)\)，故其法向量可取 \(\bs{n}=\overrightarrow{BD}\times\overrightarrow{BC_1}=(2s^2,2s^2,-s^2)\).直线 \(CD\) 的方向向量可取 \(\bs{u}=(1,0,0)\).设直线与平面的夹角为 \(\theta\)，则 \(\sin\theta=\frac{|\bs{u}\cdot\bs{n}|}{|\bs{u}|\,|\bs{n}|}=\frac{2s^2}{s^2\sqrt{4+4+1}}=\frac{2}{3}\).故选 A.
\end{solution}

\begin{problem}
已知抛物线 \(C:y^2=8x\) 与点 \(M(-2,2)\)，过 \(C\) 的焦点且斜率为 \(k\) 的直线与 \(C\) 交于 \(A\)，\(B\) 两点. 若 \(\overrightarrow{MA}\cdot\overrightarrow{MB}=0\)，则 \(k=\)（\quad）

\choices
    {\(\frac{1}{2}\)}
    {\(\frac{\sqrt{2}}{2}\)}
    {\(\sqrt{2}\)}
    {\(2\)}
\end{problem}

\begin{answer}
D
\end{answer}

\begin{solution}
由 \(y^2=8x\) 得抛物线焦点为 \(F=(2,0)\).设过焦点的直线为 \(y=k(x-2)\)，令 \(t=x-2\)，则交点对应的两个参数 \(t_1\)，\(t_2\) 满足 \(k^2t^2=8(t+2)\)，即 \(k^2t^2-8t-16=0\).因此 \(t_1+t_2=\frac{8}{k^2}\)，\(t_1t_2=-\frac{16}{k^2}\).交点坐标为 \((t_i+2,kt_i)\)，所以
\[
\begin{aligned}
\overrightarrow{MA}\cdot\overrightarrow{MB}
&=(t_1+4)(t_2+4)+(kt_1-2)(kt_2-2)\\
&=(1+k^2)t_1t_2+(4-2k)(t_1+t_2)+20\\
&=\frac{4(k-2)^2}{k^2}.
\end{aligned}
\]
直线与抛物线有两个交点，故 \(k\neq0\).由内积为零得 \(k=2\)，故选 D.
\end{solution}

\begin{problem}
已知函数 \(f(x)=\cos x\sin 2x\)，下列结论中错误的是（\quad）

\choices
    {\(y=f(x)\) 的图象关于点 \((\pi,0)\) 中心对称}
    {\(y=f(x)\) 的图象关于 \(x=\frac{\pi}{2}\) 对称}
    {\(f(x)\) 的最大值为 \(\frac{\sqrt{3}}{2}\)}
    {\(f(x)\) 既是奇函数，又是周期函数}
\end{problem}

\begin{answer}
C
\end{answer}

\begin{solution}
由 \(f(x)=2\sin x\cos^2x\)，令 \(t=\sin x\in[-1,1]\)，则 \(f(x)=2t(1-t^2)=2t-2t^3\).记 \(g(t)=2t-2t^3\)，则 \(g'(t)=2-6t^2\).在 \([-1,1]\) 上，临界点为 \(t=\pm\frac{1}{\sqrt3}\)，且 \(g(-1)=g(1)=0\)，所以最大值为 \(g\left(\frac{1}{\sqrt3}\right)=\frac{4}{3\sqrt3}=\frac{4\sqrt3}{9}\)，不是 \(\frac{\sqrt3}{2}\)，故 C 错误.

又因为 \(f(2\pi-x)=-f(x)\)，图象关于 \((\pi,0)\) 中心对称；因为 \(f(\pi-x)=f(x)\)，图象关于 \(x=\frac{\pi}{2}\) 对称；因为 \(f(-x)=-f(x)\)，所以 \(f\) 是奇函数；并且 \(f(x+2\pi)=f(x)\)，所以 \(f\) 是周期函数.因此其余结论正确，故选 C.
\end{solution}

\section{填空题}

\begin{problem}
已知 \(\alpha\) 是第三象限角，\(\sin\alpha=-\frac{1}{3}\)，则 \(\cot\alpha=\)\fillinblank{}.
\end{problem}

\begin{answer}
\(2\sqrt{2}\)
\end{answer}

\begin{solution}
因为 \(\alpha\) 在第三象限，所以 \(\cos\alpha<0\).由 \(\cos^2\alpha=1-\sin^2\alpha=1-\frac{1}{9}=\frac{8}{9}\)，得 \(\cos\alpha=-\frac{2\sqrt2}{3}\).于是 \(\cot\alpha=\frac{\cos\alpha}{\sin\alpha}=\frac{-2\sqrt2/3}{-1/3}=2\sqrt2\).
\end{solution}

\begin{problem}
\(6\) 个人排成一行，其中甲、乙两人不相邻的不同排法共有\fillinblank{} 种. （用数字作答）
\end{problem}

\begin{answer}
480
\end{answer}

\begin{solution}
\(6\) 个人任意排队有 \(6!\) 种.把甲、乙看成一个整体时有 \(5!\) 种排列，甲、乙在整体内还有 \(2!\) 种顺序，所以甲、乙相邻的排法有 \(2\times5!\) 种.因此不相邻的排法为 \(6!-2\times5!=720-240=480\) 种.
\end{solution}

\begin{problem}
记不等式组 \(\begin{cases}
x\geqslant0,\\
x+3y\geqslant4,\\
3x+y\leqslant4
\end{cases}\) 表示的平面区域为 \(D\)，若直线 \(y=a(x+1)\) 与 \(D\) 有公共点，则 \(a\) 的取值范围是\fillinblank{}.
\end{problem}

\begin{answer}
\(\left[\frac{1}{2},4\right]\)
\end{answer}

\begin{solution}
区域 \(D\) 是三角形，其顶点为 \(P=\left(0,\frac{4}{3}\right)\)，\(Q=(0,4)\)，\(R=(1,1)\).直线 \(y=a(x+1)\) 都经过定点 \(E=(-1,0)\)，所以 \(a=\frac{y}{x+1}\) 是点 \(E\) 与区域内点连线的斜率.

三个顶点对应的斜率分别为 \(\frac{4}{3}\)，\(4\)，\(\frac{1}{2}\).又因为每个区域内点都是三个顶点的凸组合，且三个顶点的横坐标都不小于 \(0\)，所以 \(\frac{y}{x+1}\) 是这三个顶点斜率的加权平均值.因此其取值范围为 \(\left[\frac{1}{2},4\right]\).
\end{solution}

\begin{problem}
已知圆 \(O\) 和圆 \(K\) 是球 \(O\) 的大圆和小圆，其公共弦长等于球 \(O\) 的半径，\(OK=\frac{3}{2}\)，且圆 \(O\) 与圆 \(K\) 所在的平面所成的一个二面角为 \(60^{\circ}\)，则球 \(O\) 的表面积等于\fillinblank{}.
\end{problem}

\begin{answer}
\(16\pi\)
\end{answer}

\begin{solution}
设球的半径为 \(R\)，两圆所在平面的交线为公共弦所在直线 \(l\).大圆半径为 \(R\)，而公共弦长为 \(R\)，故从球心 \(O\) 到 \(l\) 的距离为 \(h=\sqrt{R^2-\left(\frac{R}{2}\right)^2}=\frac{\sqrt3}{2}R\).

小圆圆心 \(K\) 是球心 \(O\) 在小圆所在平面上的射影，所以 \(OK\) 等于球心到该平面的距离.在垂直于 \(l\) 的截面内，两平面的夹角为 \(60^{\circ}\)，故 \(OK=h\sin60^{\circ}=\frac{\sqrt3}{2}R\cdot\frac{\sqrt3}{2}=\frac{3}{4}R\).由 \(OK=\frac{3}{2}\)，得 \(R=2\)，因此球的表面积为 \(4\pi R^2=16\pi\).
\end{solution}

\section{解答题}

\begin{problem}
等差数列 \(\{a_n\}\) 的前 \(n\) 项和为 \(S_n\)，已知 \(S_3=a_2^2\)，且 \(S_1\)，\(S_2\)，\(S_4\) 成等比数列，求 \(\{a_n\}\) 的通项公式.
\end{problem}

\begin{solution}
设等差数列的首项为 \(a_1\)，公差为 \(d\)，则 \(S_1=a_1\)，\(S_2=2a_1+d\)，\(S_3=3a_1+3d\)，\(S_4=4a_1+6d\).

由 \(S_1\)，\(S_2\)，\(S_4\) 成等比数列，得 \(S_2^2=S_1S_4\)，即 \((2a_1+d)^2=a_1(4a_1+6d)\)，化简得 \(d(d-2a_1)=0\).

当 \(d=0\) 时，\(a_2=a_1\)，由 \(S_3=a_2^2\) 得 \(3a_1=a_1^2\)，所以 \(a_1=0\) 或 \(a_1=3\).若 \(a_1=0\)，则 \(S_1=S_2=S_4=0\)，相邻项之比无意义，不满足等比数列的定义，故 \(a_1=3\)，从而 \(a_n=3\).

当 \(d=2a_1\) 时，\(a_2=3a_1\)，\(S_3=9a_1\).由 \(S_3=a_2^2\) 得 \(9a_1=9a_1^2\)，所以 \(a_1=0\) 或 \(a_1=1\).同样舍去全零情形，得 \(a_1=1\)，\(d=2\)，于是 \(a_n=a_1+(n-1)d=2n-1\).

因此 \(\{a_n\}\) 的通项公式为 \(a_n=3\) 或 \(a_n=2n-1\).
\end{solution}

\begin{problem}
设 \(\triangle ABC\) 的内角 \(A\)，\(B\)，\(C\) 的对边分别为 \(a\)，\(b\)，\(c\)，\((a+b+c)(a-b+c)=ac\).

\begin{enumerate}
\item 求 \(B\)；
\item 若 \(\sin A\sin C=\frac{\sqrt{3}-1}{4}\)，求 \(C\).
\end{enumerate}
\end{problem}

\begin{solution}
由余弦定理 \(b^2=a^2+c^2-2ac\cos B\)，且 \((a+b+c)(a-b+c)=(a+c)^2-b^2\).所以题设条件给出 \(a^2+2ac+c^2-b^2=ac\)，即 \(2ac(1+\cos B)=ac\).三角形三边均为正数，故 \(2(1+\cos B)=1\)，从而 \(\cos B=-\frac{1}{2}\).因为 \(0<B<\pi\)，所以 \(B=120^{\circ}\).

由 \(A+C=60^{\circ}\) 及积化和差公式，\(\sin A\sin C=\frac{\cos(A-C)-\cos(A+C)}{2}=\frac{\cos(A-C)-1/2}{2}\).代入已知条件得 \(\cos(A-C)=\frac{\sqrt3}{2}\).又因为 \(-60^{\circ}<A-C<60^{\circ}\)，故 \(A-C=\pm30^{\circ}\).结合 \(A+C=60^{\circ}\)，得到 \((A,C)=(45^{\circ},15^{\circ})\) 或 \((15^{\circ},45^{\circ})\).因此 \(C=15^{\circ}\) 或 \(C=45^{\circ}\).
\end{solution}

\begin{problem}
如图，四棱锥 \(P-ABCD\) 中，\(\angle ABC=\angle BAD=90^{\circ}\)，\(BC=2AD\)，\(\triangle PAB\) 与 \(\triangle PAD\) 都是等边三角形.

\begin{enumerate}
\item 证明：\(PB\perp CD\)；
\item 求二面角 \(A-PD-C\) 的大小.
\end{enumerate}

\FigureLayoutDeclare{img/2013/syllabus_science/q19_fig1.png}{5.8cm}{26bp 53bp 92bp 37bp}
\bitmapfigure[max width=0.42\linewidth]{img/2013/syllabus_science/q19_fig1.png}
\end{problem}

\begin{solution}
由两个等边三角形共用边 \(PA\)，得 \(AB=PA=AD\).几何形状在相似伸缩下不变，取 \(AB=AD=2\).由 \(AD\perp AB\)、\(BC\perp AB\) 以及四边形 \(ABCD\) 共面，得 \(AD\parallel BC\).建立直角坐标系，取 \(A=(0,0,0)\)，\(B=(2,0,0)\)，\(D=(0,2,0)\)，\(C=(2,4,0)\).

设 \(P=(x,y,z)\).由 \(PA=PB=PD=2\)，比较平方距离得 \(x=1\)，\(y=1\)，再由 \(PA=2\) 得 \(z=\sqrt2\)，故可取 \(P=(1,1,\sqrt2)\).

\begin{enumerate}
\item 取 \(\overrightarrow{BP}=(-1,1,\sqrt2)\)，\(\overrightarrow{CD}=(-2,-2,0)\)，则 \(\overrightarrow{BP}\cdot\overrightarrow{CD}=2-2=0\).因此 \(PB\perp CD\).
\item 设所求二面角 \(A-PD-C\) 为 \(\theta\).取 \(\bs{d}=\overrightarrow{PD}=(-1,1,-\sqrt2)\)，\(\bs{u}=\overrightarrow{PA}=(-1,-1,-\sqrt2)\)，\(\bs{v}=\overrightarrow{PC}=(1,3,-\sqrt2)\).在垂直于棱 \(PD\) 的截面内，\(\bs{u}\) 和 \(\bs{v}\) 的投影分别为
\[
\bs{u}_{\perp}=\bs{u}-\frac{\bs{u}\cdot\bs{d}}{\bs{d}\cdot\bs{d}}\bs{d}=\left(-\frac12,-\frac32,-\frac{\sqrt2}{2}\right)
\]
和
\[
\bs{v}_{\perp}=\bs{v}-\frac{\bs{v}\cdot\bs{d}}{\bs{d}\cdot\bs{d}}\bs{d}=(2,2,0)
\]
因此 \(\bs{u}_{\perp}\cdot\bs{v}_{\perp}=-4\)，\(|\bs{u}_{\perp}|=\sqrt3\)，\(|\bs{v}_{\perp}|=2\sqrt2\)，从而 \(\cos\theta=\frac{-4}{\sqrt3\cdot2\sqrt2}=-\frac{\sqrt6}{3}\).所以 \(\theta=\arccos\left(-\frac{\sqrt6}{3}\right)=\pi-\arccos\frac{\sqrt6}{3}\).
\end{enumerate}
\end{solution}

\begin{problem}
甲、乙、丙三人进行羽毛球练习赛，其中两人比赛，另一人当裁判. 每局比赛结束时，负的一方在下一局当裁判. 设各局中双方获胜的概率均为 \(\frac{1}{2}\)，各局比赛的结果都相互独立，第 1 局甲当裁判.

\begin{enumerate}
\item 求第 4 局甲当裁判的概率；
\item \(X\) 表示前 4 局中乙当裁判的次数. 求 \(X\) 的数学期望.
\end{enumerate}
\end{problem}

\begin{solution}
记第 \(i\) 局的裁判为 \(R_i\).若第 \(i\) 局的裁判确定，则另外两人比赛，下一局的裁判是负者，所以对任意不同的两人 \(U\)，\(V\)，都有 \(P(R_{i+1}=U\mid R_i=V)=\frac12\).

\begin{enumerate}
\item 已知 \(P(R_1=\text{甲})=1\)，所以 \(P(R_2=\text{乙})=P(R_2=\text{丙})=\frac12\).于是 \(P(R_3=\text{甲})=\frac12\)，且 \(P(R_3=\text{乙})=P(R_3=\text{丙})=\frac14\).因此
\[
P(R_4=\text{甲})=\frac12P(R_3=\text{乙})+\frac12P(R_3=\text{丙})=\frac14
\]
\item 由期望的线性性质，\(E(X)=\sum_{i=1}^4P(R_i=\text{乙})\).其中 \(P(R_1=\text{乙})=0\)，\(P(R_2=\text{乙})=\frac12\)，\(P(R_3=\text{乙})=\frac14\)，以及
\[
P(R_4=\text{乙})=\frac12P(R_3=\text{甲})+\frac12P(R_3=\text{丙})=\frac12\times\frac12+\frac12\times\frac14=\frac38
\]
所以 \(E(X)=0+\frac12+\frac14+\frac38=\frac98\).
\end{enumerate}
\end{solution}

\begin{problem}
已知双曲线 \(C:\frac{x^2}{a^2}-\frac{y^2}{b^2}=1\)（\(a>0\)，\(b>0\)）的左、右焦点分别为 \(F_1\)，\(F_2\)，离心率为 3，直线 \(y=2\) 与 \(C\) 的两个交点间的距离为 \(\sqrt{6}\).

\begin{enumerate}
\item 求 \(a\)，\(b\)；
\item 设过 \(F_2\) 的直线 \(l\) 与 \(C\) 的左、右两支分别交于 \(A\)，\(B\) 两点，且 \(|AF_1|=|BF_1|\)，证明：\(|AF_2|\)，\(|AB|\)，\(|BF_2|\) 成等比数列.
\end{enumerate}
\end{problem}

\begin{solution}
\begin{enumerate}
\item 设双曲线的焦距参数为 \(c\)，则 \(c^2=a^2+b^2\)，\(\frac{c}{a}=3\).所以 \(c=3a\)，\(b^2=8a^2\).将 \(y=2\) 代入双曲线方程，得 \(\frac{x^2}{a^2}-\frac{4}{b^2}=1\)，即 \(x^2=a^2\left(1+\frac{4}{b^2}\right)\).两交点的横坐标相反，故其距离为 \(2a\sqrt{1+\frac{4}{b^2}}=\sqrt6\).代入 \(b^2=8a^2\)，得 \(4a^2+2=6\)，所以 \(a=1\)，\(b=2\sqrt2\).
\item 由第一问，\(F_1=(-3,0)\)，\(F_2=(3,0)\).设直线 \(l\) 为 \(y=k(x-3)\)，并记它与双曲线交点的横坐标为 \(x_A\)，\(x_B\).代入双曲线方程，得
\[
(8-k^2)x^2+6k^2x-(9k^2+8)=0
\]
因为 \(A\)，\(B\) 在双曲线左右两支上，所以 \(x_A\ne x_B\).又由 \(y^2=8(x^2-1)\)，点 \((x,y)\) 到 \(F_1\) 的距离平方为 \((x+3)^2+y^2=9x^2+6x+1\).条件 \(|AF_1|=|BF_1|\) 给出 \(9x_A^2+6x_A+1=9x_B^2+6x_B+1\)，所以 \(x_A+x_B=-\frac23\).

由韦达定理，\(-\frac{6k^2}{8-k^2}=-\frac23\)，解得 \(k^2=\frac45\).进一步 \(x_Ax_B=-\frac{9k^2+8}{8-k^2}=-\frac{19}{9}\)，从而 \((x_A-x_B)^2=(x_A+x_B)^2-4x_Ax_B=\frac{80}{9}\).于是可取 \(x_A=\frac{-1-2\sqrt5}{3}\)，\(x_B=\frac{-1+2\sqrt5}{3}\)，故 \(A\)，\(B\) 均在 \(F_2\) 的左侧，且 \(B\) 位于 \(A\) 与 \(F_2\) 之间.

直线 \(l\) 的方向因子为 \(\sqrt{1+k^2}=\frac{3}{\sqrt5}\)，所以 \(|AB|=\sqrt{1+k^2}|x_A-x_B|=4\).又因为 \(3-x_A>0\)，\(3-x_B>0\)，故
\[
\begin{aligned}
|AF_2|\,|BF_2|
&=(1+k^2)(3-x_A)(3-x_B)\\
&=\frac95\left(9-3(x_A+x_B)+x_Ax_B\right)\\
&=\frac95\left(9+2-\frac{19}{9}\right)=16.
\end{aligned}
\]
因此 \(|AF_2|\,|BF_2|=16=|AB|^2\)，即 \(|AF_2|\)，\(|AB|\)，\(|BF_2|\) 成等比数列.
\end{enumerate}
\end{solution}

\section{选考题：解答题}

\begin{problem}
已知函数 \(f(x)=\ln\left(1+x\right)-\frac{x(1+\lambda x)}{1+x}\).

\begin{enumerate}
\item 若 \(x\geqslant0\) 时 \(f(x)\leqslant0\)，求 \(\lambda\) 的最小值；
\item 设数列 \(\{a_n\}\) 的通项 \(a_n=1+\frac{1}{2}+\frac{1}{3}+\dots+\frac{1}{n}\)，证明：\(a_{2n}-a_n+\frac{1}{4n}>\ln2\).
\end{enumerate}
\end{problem}

\begin{solution}
\begin{enumerate}
\item 在 \(x\geqslant0\) 上，
\[
\begin{aligned}
f'(x)
&=\frac{1}{1+x}-\frac{1+2\lambda x+\lambda x^2}{(1+x)^2}\\
&=\frac{x(1-2\lambda-\lambda x)}{(1+x)^2}.
\end{aligned}
\]
又 \(f(0)=0\).若 \(\lambda\geqslant\frac12\)，则对 \(x\geqslant0\)，有 \(1-2\lambda-\lambda x\leqslant0\)，从而 \(f'(x)\leqslant0\)，故 \(f(x)\leqslant f(0)=0\).

若 \(\lambda<\frac12\)，则 \(1-2\lambda>0\)，所以当 \(x>0\) 足够接近于 \(0\) 时，\(1-2\lambda-\lambda x>0\)，从而 \(f'(x)>0\)，于是 \(f(x)>f(0)=0\)，与题设矛盾.因此 \(\lambda\geqslant\frac12\)，最小值为 \(\frac12\).
\item 因为函数 \(g(x)=\frac{1}{x}\) 在 \((0,+\infty)\) 上严格凸，所以对每个正整数 \(k\)，都有
\[
\int_k^{k+1}\frac{1}{x}\,\mathrm{d}x<\frac12\left(\frac{1}{k}+\frac{1}{k+1}\right)
\]
对 \(k=n\)，\(n+1\)，\(\dots\)，\(2n-1\) 的不等式求和，得
\[
\begin{aligned}
\ln2
&=\int_n^{2n}\frac{1}{x}\,\mathrm{d}x\\
&<\frac{1}{2n}+\sum_{k=n+1}^{2n-1}\frac{1}{k}+\frac{1}{4n}\\
&=\sum_{k=n+1}^{2n}\frac{1}{k}+\frac{1}{4n}\\
&=a_{2n}-a_n+\frac{1}{4n}.
\end{aligned}
\]
故不等式得证.
\end{enumerate}
\end{solution}
